\subsection{Presentation de la problématique}
\begin{frame}{Problématique}
    \begin{alertblock}{D'après le ministère de la Santé : }
        \begin{itemize}
            \item \textbf{31 millions} d'appels d'urgence en 2018
            \item Seuls \textbf{69\%} des appels décrochés dans la minute
        \end{itemize}
    \end{alertblock}
    \begin{block}{Objectif}
        Utiliser la reconnaissance vocale par réseau de neurones pour aider à classifier rapidement l'objet d'un appel.
    \end{block}
    \openup 2em
    \begin{center}
        \centering
        \includegraphics[width=250px]{4-phone.png}
    \end{center}


\end{frame}


\begin{frame}{Données}

	\begin{block}{12 Mots-clés}
		Le site du gouvernement recense ces mots-clés pour les numéros d'appel d'urgence : \\
		\begin{description}
			\item[n°15] malaise, hémorragie, brûlure, intoxication
			\item[n°17] violences, agression, vol, cambriolage
			\item[n°18] incendie, gaz, effondrement, électrocution
		\end{description}
	\end{block}

\end{frame}